% !TEX root = ../main.tex
% ---------------------------------------------------------------------
%  B2 -- SPS und TwinCAT (Skript Kap. 2, Anhang B, D, F; AS 1)
% ---------------------------------------------------------------------
\vlsec{cB2}{B2}{SPS und TwinCAT}
\label{sec:twincat}

\begin{formel}[Scan-Zyklus einer SPS]
\T{1.} Alle Eingänge ins \T{Prozessabbild der Eingänge} (PAE) lesen
\ra \T{2.} Programm auf diesem eingefrorenen Abbild abarbeiten
\ra \T{3.} \T{Prozessabbild der Ausgänge} (PAA) auf die Klemmen schreiben.
Wiederholt mit fester Zykluszeit. Eine Eingangsänderung \emph{während}
des Zyklus wird erst im nächsten Zyklus sichtbar.
\end{formel}

\SW{TwinCAT} Soft-SPS: Echtzeitkern \emph{unter} Windows. Engineering
(Visual Studio + TwinCAT XAE) auf dem Entwicklungsrechner (ER), Runtime auf
dem Zielsystem (Labor: CX5010/CX9020). TwinCAT ist \emph{Software}, kein
Feldbus -- der Feldbus heisst EtherCAT.

\SW{LEDs} PWR grün $=$ Strom. TC \T{grün} $=$ Run-Modus (Boot-Projekt
läuft), TC \T{blau} $=$ Config-Modus.
\SW{Modi} \T{Config} (mit \T{Freerun}: Klemmen ohne Programm testen,
Reiter \emph{Online}), \T{Run} (Programm läuft).

\subsection{Vom leeren Projekt zum laufenden Programm}
\begin{rezept}[Ablauf (AS 1)]
\begin{enumerate}
\item ER und SPS ins \T{gleiche Subnetz} (statische IP), \R{ping}.
\item TwinCAT-XAE-Projekt \ra Zielsystem wählen \ra suchen \ra
  \T{Route zufügen} (siehe ADS unten).
\item E/A \ra Gerät 1 (EtherCAT) \ra \T{Scan} / \emph{Nach Boxen suchen}.
\item SPS-Projekt hinzufügen (\R{MAIN} entsteht), Task-Zykluszeit setzen
  (z.\,B.\ 5 Zyklusticks $=$ 5\,ms).
\item GVL mit \R{AT \%I*} / \R{AT \%Q*}-Variablen, Code schreiben.
\item \T{Projektmappe erstellen} (kompilieren).
\item \emph{Instance} \ra PlcTask Inputs/Outputs \ra \T{Verknüpfung ändern}
  -- es werden nur typ- und grössengleiche Prozessdaten angeboten.
\item \T{Konfiguration aktivieren} (Zuordnungen), Run-Modus.
\item SPS \ra \T{Einloggen} (lädt Programm) \ra \T{Start}.
\end{enumerate}
Jede Codeänderung: neu erstellen und neu laden. Online kann man Werte
\T{schreiben} oder \T{forcieren} (testen ohne Hardware).
\end{rezept}

\subsection{E/A-Variablen und Verknüpfung}
\label{rez:gvlio}
\begin{lstlisting}
{attribute 'qualified_only'}  // Zugriff nur als GVL_IO.xxx
VAR_GLOBAL
  bTasterGelb   AT %I* : BOOL;  // Eingang (EL1004)
  bLampeRot     AT %Q* : BOOL;  // Ausgang (EL2004)
  iSensorRoh    AT %I* : INT;   // Analog-Eingang (EL3004)
  iAnalogOut    AT %Q* : INT;   // Analog-Ausgang (EL4034/EL3034)
END_VAR
VAR_GLOBAL CONSTANT
  rT : REAL := 0.01;            // Abtastzeit (s)
END_VAR
\end{lstlisting}
\begin{formel}[AT-Adressen]
\R{\%I*} $=$ Eingang, \R{\%Q*} $=$ Ausgang, \R{*} $=$ Adresse vergibt
TwinCAT. Die Zuordnung zur Klemme passiert im Engineering, nicht im Code
\ra gleiche Software auf anderer Hardware. Typ muss passen: 16-Bit-Analog
\ra \R{INT}, Digital \ra \R{BOOL}. Ohne \R{AT} erscheint die Variable
\emph{nicht} im Prozessabbild.
\end{formel}
\begin{falle}[GVL ist nicht nur für Klemmen]
Globale Variablenlisten enthalten auch Konstanten
(\R{VAR\_GLOBAL CONSTANT}) und normale globale Daten. Mit
\R{qualified\_only} muss man \R{GVL\_IO.} immer davor schreiben.
\end{falle}

\SW{Analoge Klemmen} 0--10\,V $\leftrightarrow$ \R{INT} 0--32767.
10\,V am Ausgang \ra Wert 32767 ($=2^{15}-1$, grösster INT).
EL3004: 4 Kanäle, 12 Bit Auflösung, Rohwert ändert sich in 16er-Schritten.

\subsection{Projektstruktur (TwinCAT XAE)}
\begin{ztab}{@{}lX@{}}
SYSTEM & Tasks, Echtzeit, Lizenzen, Routen \\
MOTION & NC-Achsen (nicht im Labor) \\
SPS & IEC-61131-Projekte \\
SAFETY & TwinSAFE-Logik (eigene Logic-Klemme) \\
E/A & EtherCAT-Klemmen, Prozessabbild \\ \midrule
References & Bibliotheken (\R{Tc2\_Standard}: \R{TON}, \R{R\_TRIG}) \\
DUTs & eigene Typen: STRUCT, ENUM \\
GVLs / POUs & globale Variablen / Programmcode \\
VISUs & Visualisierung (+ VISU\_TASK, eigener langsamer Task) \\
.tmc & Symbolbeschreibung für ADS-Zugriff von aussen \\
Instance & mit Hardware verknüpfte Instanz (PlcTask In/Out) \\
\end{ztab}

\SW{ADS} Beckhoff-Protokoll über TCP/IP für Download, Start/Stopp,
Online-Monitoring. Ziel $=$ (\T{AMS Net Id}, ADS-Port); AMS Net Id
z.\,B.\ \R{5.62.71.101.1.1} (oft IP + \R{.1.1}), bleibt bei IP-Wechsel.
\T{Route} $=$ Vertrauensbeziehung, sonst nimmt die SPS keine ADS-Befehle an.
Der \T{TwinCAT-Systemdienst} läuft immer (AMS-Router), auch ohne Visual
Studio.

\SW{CoE} Objektverzeichnis der Klemme (Index:Subindex, z.\,B.\
\R{0x8000:02}). EL3004 \R{0x80n0}: \R{:02} Presentation (\T{Signed} $=$
Zweierkomplement $=$ INT, darum direkt \R{INT\_TO\_REAL}), \R{:06}/\R{:15}
Filter, \R{:07} Grenzwert \ra Statusbit (die Klemme erkennt das
\emph{selbst}, ohne SPS). CoE-Online wirkt sofort, aber nur bis Neustart;
dauerhaft \ra \T{CoE-Startup}-Liste.

\subsection{Remanente Variablen}
\label{sec:remanenz}
\begin{ztab}{@{}lccc@{}}
Ereignis & normal & \R{RETAIN} & \R{PERSISTENT} \\ \midrule
Online Change & bleibt & bleibt & bleibt \\
Reset warm & weg & bleibt & bleibt \\
Reset kalt & weg & weg & bleibt \\
Neuer Download & weg & weg & bleibt \\
Reset Ursprung & weg & weg & weg \\
\end{ztab}
\begin{lstlisting}
VAR RETAIN     rIntegrator  : REAL; END_VAR
VAR PERSISTENT diStueckzahl : DINT; END_VAR
\end{lstlisting}
\SW{Wann was} Stückzähler über Programm-Update \ra \R{PERSISTENT}.
I-Anteil eines Reglers über kurzen Neustart \ra \R{RETAIN}.
\T{USV/UPS}: Akku/Kondensator überbrückt Stromausfall, Zeit zum Sichern.
\T{NOVRAM}: nichtflüchtiger RAM, jeden Zyklus beschreibbar.
